% HD-G29. GPT-6 Astra Pro /root/g34; inherited effort; 21 November 2026.
\section{HD-G29: actual field descent of the two-dimensional blocks}
\label{sec:two-block-field-descent}

In the complete setting of HD-G28, every original complex
two-dimensional block is the scalar extension of an actual
two-dimensional \(K\)-subspace. Every actual \(K\)-linear
overgroup normalizing \(N\) permutes these \(K\)-blocks.
Consequently, for every imaginary quadratic \(K\), a natural
primitive sixfold sphere-cover source whose full coefficient
group normalizes such an absolutely irreducible \(N\) has
period rank at most seven. This holds for every \(r\ge5\),
without an upper bound on the coefficient rank \(2r\).

The argument uses the characteristic kernel \(T\) and the
literal tree coordinates of HD-G28. The field descent is proved
from actual matrix projectors. It is not an assertion that
arbitrary complex block systems are Galois stable.

\paragraph{1. The exact list of irreducible kernel modules.}
In the tree coordinates, write
\[
   V_{\mathbf C}=\bigoplus_{i=1}^rV_i,\qquad
   T\subset L^r,\qquad (L')^r\subset T,\qquad N=T\rtimes S_r.
                                                               \tag{G29.1}
\]
All these are the actual maps and inclusions proved in HD-G28.
The common finite coordinate image \(L\subset GL_2(\mathbf C)\)
is irreducible. Each \(V_i\) is therefore an irreducible
\(T\)-module: its coordinate image is all of \(L\).
The finite group \(L\) is not abelian, since commuting
finite-order complex matrices are simultaneously diagonalizable,
which would give an invariant line. Thus \(L'\ne1\).

Choose \(a\in L'\setminus\{1\}\). The inclusion in
(G29.1) supplies the actual element
\[
             a_i=\operatorname{diag}(1,\ldots,1,a,1,\ldots,1)
                         \in T.                            \tag{G29.2}
\]
It acts nontrivially on \(V_i\) and as the identity on
\(V_j\) for \(j\ne i\). An intertwiner between \(V_i\)
and \(V_j\) would make these two actions conjugate, which
is impossible for \(a\ne1\). Hence the \(r\) modules
\(V_i\) are pairwise inequivalent. They are exactly the
irreducible constituents of the specified \(T\)-module
\(V_{\mathbf C}\), each with multiplicity one.

\paragraph{2. Literal projectors and complexification of the original algebra.}
Let \(\rho_i\) be the actual representation on \(V_i\),
and \(\chi_i(t)=\operatorname{tr}\rho_i(t)\). Define
\[
       p_i=\frac{2}{|T|}\sum_{t\in T}
                            \chi_i(t^{-1})\rho(t).          \tag{G29.3}
\]
This is precisely the projection onto \(V_i\) along all
the other blocks. Here is the exact verification. The coefficient
function is constant on conjugacy classes, so the restriction
of \(p_i\) to each irreducible \(V_j\) is scalar. Averaging
\(\operatorname{Hom}(V_i,V_j)\) by
\[
           A\longmapsto |T|^{-1}\sum_{t\in T}
                              \rho_j(t)A\rho_i(t^{-1})
                                                               \tag{G29.4}
\]
is the idempotent onto the intertwiner space. Its trace is
\(|T|^{-1}\sum_t\chi_j(t)\chi_i(t^{-1})=\delta_{ij}\),
by irreducibility and the inequivalence just proved. Thus
\(\operatorname{tr}(p_i|_{V_j})=2\delta_{ij}\), and the
scalar is \(\delta_{ij}\). In particular
\[
       p_i p_j=\delta_{ij}p_i,\qquad
       \sum_i p_i=I,\qquad \operatorname{rank}p_i=2.        \tag{G29.5}
\]

For clarity, these are exactly the primitive central idempotents
of the image algebra of \(T\), with no omitted additional
kernel types. The literal restriction map identifies
\[
 \mathbf C\rho(T)=\bigoplus_{i=1}^r
                         \operatorname{End}_{\mathbf C}(V_i).
                                                               \tag{G29.6}
\]
Surjectivity can be checked by matrix units, rather than inferred
from a dimension count. Choose bases of the \(V_i\); averaging
as in (G29.4) gives
\[
 E^{(i)}_{ab}=\frac{2}{|T|}\sum_{t\in T}
                  (\rho_i(t^{-1}))_{ba}\rho(t),
                      \qquad 1\le a,b\le2.                 \tag{G29.7}
\]
On \(V_i\) this is the ordinary matrix unit \(E_{ab}\),
and on every \(V_j\), \(j\ne i\), it is zero. Thus
the displayed restriction map has the stated explicit inverse
on every matrix-unit basis.

Put \(\mathcal A=K\rho(T)\subset M_{2r}(K)\).
The natural scalar-extension map
\(\mathcal A\otimes_K\mathbf C\to M_{2r}(\mathbf C)\)
is injective: a \(K\)-basis of a matrix subspace remains
linearly independent after extending the field, as follows
from any nonzero maximal minor of its coefficient matrix.
Its image is exactly (G29.6). This is the complexification
of the original \(K\)-algebra, not a separately chosen
complex endomorphism algebra.

\paragraph{3. The exact Galois action fixes every projector.}
Let \(h\) be the exponent of \(T\) and let
\(E=K(\zeta_h)\subset\mathbf C\). It is a finite Galois
extension of \(K\), since it is a splitting field for
\(X^h-1\). Each \(\chi_i(t)\) is a sum of \(h\)-th
roots of unity, so (G29.3) gives
\(p_i\in\mathcal A_E=E\rho(T)\).
The centre of \(\mathcal A_E\) is exactly
\[
                         Z(\mathcal A_E)=\bigoplus_i E p_i.
                                                               \tag{G29.8}
\]
Indeed the centre is the kernel of the finite system of
commutator equations with a basis of \(\mathcal A\).
Their rank is preserved by scalar extension. Equation
(G29.6) makes its dimension \(r\), and the \(r\)
independent idempotents in (G29.5) give the displayed basis.
They are its complete list of primitive central idempotents.

For \(\gamma\in\operatorname{Gal}(E/K)\), applying
\(\gamma\) to matrix entries fixes each original
\(\rho(t)\), and defines a semilinear automorphism of
\(\mathcal A_E\). It therefore permutes exactly the
\(r\) idempotents in (G29.8). Denote this permutation
by \(\pi_\gamma\). Conjugation by each original
\(\rho(n)\), \(n\in N\), also permutes these idempotents
according to its block permutation. Since \(\rho(n)\)
has entries in \(K\), these two actions commute literally:
\[
 \gamma\bigl(\rho(n)p_i\rho(n)^{-1}\bigr)
       =\rho(n)\gamma(p_i)\rho(n)^{-1}.                    \tag{G29.9}
\]
The block permutations of \(N\) are the full natural
\(S_r\). Its centralizer in the permutation group on
these \(r\) labels is trivial. To see this explicitly,
the stabilizer of label \(i\) is \(S_{r-1}\), whose
only fixed label is \(i\) for \(r\ge3\). A commuting
permutation must send \(i\) to a label fixed by that
stabilizer, so must fix \(i\); this holds for all \(i\).
Hence \(\pi_\gamma=1\) for every \(\gamma\), and
\[
            p_i\in M_{2r}(K),\qquad
            U_i=p_iK^{2r},\qquad
            \dim_KU_i=2,\qquad U_i\otimes_K\mathbf C=V_i.
                                                               \tag{G29.10}
\]
The last equality follows from the actual projector map,
not just equality of dimensions. Also
\(\chi_i(t)=\operatorname{tr}(p_i\rho(t))\in K\).

\paragraph{4. Every actual normalizer preserves the descended blocks.}
HD-G28 proves \(T\) characteristic in \(N\).
Thus an actual \(g\in GL_{2r}(K)\) normalizing \(N\)
normalizes \(T\) and the original algebra \(\mathcal A\).
Its conjugation permutes the complete primitive central
idempotents of \(\mathcal A\), which are precisely the
\(p_i\) now in (G29.10). Therefore \(gU_i=U_j\) for
the corresponding permutation. This is a transitive block
system for every such overgroup containing \(N\), since
\(N\) already acts as \(S_r\).

For a finite coefficient overgroup \(G\), choose
\(U=U_1\), let \(H\) be its stabilizer, and let
\(\sigma:H\to GL_K(U)\) be the restriction. Its
\(T\)-submodule is irreducible, so \(\sigma\) is
absolutely irreducible. For actual left coset representatives
\(t_j\) the map
\[
        \bigoplus_{j=1}^rU\longrightarrow V_K,
                  (u_j)\longmapsto\sum_j\rho(t_j)u_j
                                                               \tag{G29.11}
\]
is an isomorphism, with block transition \(\sigma(h)\)
when \(gt_j=t_kh\). Thus this is literally
\(V_K=\operatorname{Ind}_H^G\sigma\).
For an imaginary quadratic exact character field \(K\),
the induced trace formula and the actual \(K\)-matrices
give \(\mathbf Q(\operatorname{tr}\sigma)=K\).

\paragraph{5. The actual quotient factor and the period-rank bound.}
Let \(\widetilde C\to\mathbf P^1\) be the original
full-\(G\) cover. With \(H_0=\ker\sigma\), take
the actual quotients
\[
       Y=\widetilde C/H,\qquad C'=\widetilde C/H_0,
                         \alpha:\widetilde C\to C'.        \tag{G29.12}
\]
The coefficient functions and inverse are exactly
\[
 \begin{aligned}
 F_v(x)&=\pi_U\rho(x)v,&
 F_v(hx)&=\sigma(h)F_v(x),\\
 F_{\rho(g)v}(x)&=F_v(xg),&
 F&\longmapsto\sum_{H\backslash G}\rho(x)^{-1}F(x).
 \end{aligned}                                              \tag{G29.13}
\]
The left-deck sheet \(Ht_j^{-1}\) corresponds to the block
\(t_jH\). These are the explicit induction and
middle-extension maps of HD-G27, paragraphs1--2; a branch
cycle contributes its product-monodromy fixed space on both
sides, so the maps include the ramified stalks.

The prescribed positive invariant Hermitian coefficient form
makes these blocks orthogonal: between two inequivalent
\(T\)-modules its cross term would give an intertwiner,
which is zero. For \(v\in U(K)\setminus\{0\}\) and
\(q=h(v,v)\), the actual rational rank-one corner
idempotents consequently satisfy
\[
 \frac{2(2r)}{|G|q}\sum_{g\in G}
       \operatorname{Re}h(v,\rho(g^{-1})v)g
   =\frac{2\cdot2}{|H|q}\sum_{h\in H}
       \operatorname{Re}h(v,\sigma(h^{-1})v)h,              \tag{G29.14}
\]
because the coefficients outside \(H\) vanish and
\(|G|=r|H|\). This is equality in \(\mathbf Q[G]\).
The same integral multiple therefore gives the same actual
abelian image inside \(J\widetilde C\). Pullback by
\(\alpha\) carries the corresponding primitive factor
of \(JC'\) isogenously onto it, with
\[
       \alpha_*\alpha^*=[|H_0|],\qquad
       (\alpha^*)^*\theta_{J\widetilde C}
                                      =|H_0|\theta_{JC'}.
                                                               \tag{G29.15}
\]
The two real Schur constants agree, since
\(2|G|/((2r)q)=2|H|/(2q)\), and the same corner
field action is retained. These exact paired Hodge maps
identify the period maps after finite base change fixing
the finite monodromy labels.

An originally specified primitive coefficient line need not
lie in \(U\). Its change to the line \(Kv\) is also an
actual polarized isogeny. For a nonzero vector \(w\) on
that original \(K\)-line, put \(q_w=h(w,w)\) and
\(T_{v,w}(x)=v\,h(w,x)/q_w\). In the rational block the
literal element
\[
 b=\frac{2(2r)}{|G|}\sum_{g\in G}
       \operatorname{Re}\operatorname{tr}
               (T_{v,w}\rho(g^{-1}))\,g                    \tag{G29.15a}
\]
has actions \(T_{v,w}\) and its conjugate on the two
coefficient representations, and zero on the other rational
blocks, by the same matrix-coefficient averaging.
Writing \(e_w,e_v\) for their orthogonal corner projectors,
one has \(b=e_vbe_w\) and
\(b^\dagger b=(q/q_w)e_w\).
An integer \(M>0\) clearing the coefficients gives the
actual Jacobian homomorphism \(a_{Mb}\). On the two
primitive factors it is an isogeny, since the displayed
coefficient-line map is invertible. The Rosati identity gives
its exact pullback polarization multiplier \(M^2q/q_w\).
Thus choosing \(v\in U\) retains the period rank of the
originally specified corner as well.

The faithful coefficient system on \(Y\) has rank two
and first parabolic cohomology of dimension six over \(K\).
Let \(g=g(Y)\), and let \(b\) be its number of nontrivial
coefficient branch values, with codimensions \(c_j\ge1\).
Irreducibility kills global invariants and dual invariants,
so the exact Euler formula is
\[
                       6=2(2g-2)+\sum_{j=1}^b c_j.         \tag{G29.16}
\]
For \(g=0\), it gives \(b\le10\), so the pointed-base
parameter dimension is at most \(b-3\le7\). For \(g=1\),
it gives \(b\le6\), and dimension at most six, including
the unpointed dimension-one case. For \(g=2\), it gives
\(b\le2\), and dimension at most \(3+b\le5\).
For \(g\ge3\), the right side is already at least eight,
which is impossible. Finite monodromy covers have discrete
choices over the fixed pointed base, so the same bounds apply
to the faithful cover family and its primitive period map.

Consequently the original period rank is at most seven.
This proves the asserted all-rank normalizer restriction
through the actual coefficient, cover, corner and
polarization maps. It does not assert that the normal subgroup
generated by bireflections is absolutely irreducible in every
larger-rank source; that remains the explicit hypothesis
in this result and is not obtained merely by taking a
normal closure.
